\documentclass[12pt, dutch]{article}
\usepackage{hyperref}
\usepackage{graphicx}
\hypersetup{
    colorlinks=true,
    linkcolor=black,
    filecolor=magenta,      
    urlcolor=blue,
    pdfpagemode=FullScreen,
    }
\urlstyle{same}
\title{SORT4: een korte analyse van Quick Sort}
\author{Manoah Kinnaert}
\graphicspath{{charts/SORT4}}
\begin{document}

\begin{titlepage}
    \maketitle
\end{titlepage}
\pagebreak
\section{Korte inleiding}
In dit (korte) verslag bespreken we een klein experimentje rond het sorteren van willekeurige reële getallen met behulp van Quick Sort. Hierbij tellen we het aantal vergelijkingen tussen elementen alsook het aantal verwisselingen
van elementen in de array. We onderzoeken hoe het aantal vergelijkingen en het aantal verwisselingen verandert in functie van de invoergrootte N (zijnde het aantal elementen van de array).
We doen dit in de hoop een idee te krijgen van hoe Quick Sort zich zal gedragen naarmate de invoergrootte steeds groter wordt. Daarnaast ook welke invloed de sturctuur van de data (of de lijst al dan niet gesorteerd is) heeft op het sorteren.
Verder kijken we welke invloeden de twee manieren van partitioneren, Hoare en Lomuto, hebben op de tijdscomplexiteit.
Alle code die gebruikt is voor dit experiment is hier terug te vinden: \\\ 
\url{https://github.com/ManoahKinnaert/Sorts}.

\section{Invoergrootte en arrays}
Voor de invoergrootte van de arrays is er gekozen voor arrays startende met lengte 10 tot en met lengte 1000. 
Een test omvat het meerdermaals uitvoeren van quicksort op een array waarbij we telkens de lengte verhogen met 10 elementen.

\section{Keuze van grafiek / weergave}
Voor dit experiment is er gekozen voor een grafiek met logaritmisch geijkte assen. Waarom logaritmisch?
Er is gekozen voor logaritmisch geijkte assen omdat het verwachte gedrag logaritmisch is. Het leek dan ook logisch om een logaritmische schaal te gebruiken om de grafiek wat netter weer te geven.

\section{Bevindingen}

\section{Bevindingen vs theorie}

\end{document}